\section{本讲小结}

极化问题需要区分自由电荷、束缚电荷和总电场。$\bm D$ 的高斯定理对应自由电荷，界面法向 $D$ 的跳变由自由面电荷决定，切向 $E$ 连续。界面方向决定分区是按面积并联还是按厚度串联。

插入介质前先判断电源连接状态。孤立电容保持电量，恒压电源保持电压；两种约束下介质都可被吸入，但场能变化的符号不同。求受力和机械功应计入电源功，不能只看场能是否下降。

非均匀介质、分层耐压与部分金属插片都可回到相同的场—势—电荷关系。使用近似时需说明边缘、线性响应、刚性与无损条件，并用极限和总束缚电荷检查结果。

